% ============================================================
% SECTION 23 — NUMERICAL FRAMEWORK
% ============================================================
\section{Numerical framework}
\label{part-5:sec:numerical-framework}

We describe the numerical methods used to obtain the results reported 
in this part. The methods are those of Part~II, Appendix~B, 
specialised to the family of truncations $K_{\max}\in\{2,3,4,5\}$.

\subsection{The truncation family}
\label{part-5:ssec:truncation-family}

For each truncation parameter $K_{\max}\in\N$, the mode set is
\begin{equation}\label{part-5:eq:mode-set-V}
\KK_{K_{\max}} := \bigl\{(k,\sigma)\in\Z^3_\ast\times\{+,-\} : 
\abs{k}\leq K_{\max}\bigr\}.
\end{equation}
The total number of modes and the real dimension of the state space 
are reported in Table~\ref{part-5:tab:truncation-params}.

\begin{table}[h]
\centering
\caption{Truncation parameters. The number of nonzero entries in the 
triadic tensor is reported in the last column.}
\label{part-5:tab:truncation-params}
\begin{tabular}{ccccc}
\toprule
$K_{\max}$ & $\#\Z^3_\ast\cap\{\abs{k}\leq K_{\max}\}$ & 
$2\abs{\KK_{K_{\max}}}$ & $\dim_\R$ & $\operatorname{nnz}(\Tt)$ \\
\midrule
$2$ & $124$ & $248$ & $248$ & $7.0\times 10^4$ \\
$3$ & $342$ & $684$ & $684$ & $6.2\times 10^5$ \\
$4$ & $728$ & $1456$ & $1456$ & $3.0\times 10^6$ \\
$5$ & $1330$ & $2660$ & $2660$ & $5.3\times 10^6$ \\
\bottomrule
\end{tabular}
\end{table}

The triadic tensor $\Tt_{\gamma\alpha\beta}$ is computed in the 
helicoidal polarisation basis by direct evaluation of the polarisation 
products, as described in Part~I, Section~3. The threshold for 
numerical support is $10^{-15}$; all tensor entries below this 
threshold are discarded.

\subsection{The optimal-state construction}
\label{part-5:ssec:optimal-construction}

The optimal states $U^{\ast,(K)}$ are computed by Riemannian gradient 
descent on the unit sphere $\AA_s$, as described in Part~II, 
Section~8. The gradient is
\begin{equation}\label{part-5:eq:grad-J}
\nabla_{\AA_s}J(U) = \Pi_{T_U\AA_s}\bigl[\Lambda^{-2s}\bigl(
\Lambda^{2s}N(U,U) + L_U^\dagger(\Lambda^{2s}U)\bigr)\bigr],
\end{equation}
where $\Pi_{T_U\AA_s}$ denotes the orthogonal projection onto the 
tangent space \eqref{part-5:eq:tangent} of the sphere, and $L_U$ is the 
linearised operator of Definition~7 of Part~I. The step size is 
chosen adaptively by Armijo backtracking; convergence is monitored 
by the residual $\norm{\nabla_{\AA_s}J(U)}_{H^s}<10^{-7}$.

\subsection{The time integration}
\label{part-5:ssec:time-integration}

The evolution of the state under the truncated Euler equations is 
computed with a fourth-order Runge--Kutta scheme with fixed step 
$\dd t = 5\times 10^{-4}$. The conservation of the $L^2$ norm is 
monitored at each step; the relative error $\abs{\Delta E/E_0}$ 
remains below $10^{-14}$ for all reported trajectories.

The coherence functional $\chi(t)$, the efficiency $\kappa(t)$, and 
the growth rate $\gamma(t)$ are evaluated at each integration step 
and stored for post-processing.

\subsection{The channel decomposition}
\label{part-5:ssec:channel-decomposition}

For each state $U(t)$, the nonlinear term $N(U)$ is decomposed into 
the pure and acopling contributions:
\begin{equation}\label{part-5:eq:decomp-N}
N(U) = N^{\mathrm{pure}}(U) + N^{\mathrm{ac}}(U),
\end{equation}
with $N^{\mathrm{pure}}$ the sum of the two pure channels $+++$ and 
$---$, and $N^{\mathrm{ac}}$ the sum of the six acopling channels. The 
corresponding contributions to the coherent channel are
\begin{equation}\label{part-5:eq:S-pure-ac}
\SS^{\mathrm{pure}} := \mathrm{Re}\ip{\Lambda^s U}{\Lambda^s N^{\mathrm{pure}}(U)},
\qquad
\SS^{\mathrm{ac}} := \mathrm{Re}\ip{\Lambda^s U}{\Lambda^s N^{\mathrm{ac}}(U)}.
\end{equation}
By Theorem~\ref{part-5:thm:decoupling} of Part~I, both vanish identically 
when $U$ is real and satisfies the reality condition.

% ============================================================
% SECTION 24 — THE OPTIMAL FUNCTIONAL
% ============================================================
\section{The optimal functional $J^\ast(K)$}
\label{part-5:sec:J-numeric}

We report the numerical values of the maximum of the coherent 
channel $J^\ast(K)$ in the truncations $K_{\max}\in\{2,3,4,5\}$.

\subsection{Definition and reference}
\label{part-5:ssec:J-def-V}

Recall from Part~I, Definition~17, the coherent channel
\begin{equation}\label{part-5:eq:J-def-V}
J(U) := \mathrm{Re}\ip{\Lambda^s U}{\Lambda^s N(U)}.
\end{equation}
The optimal value is
\begin{equation}\label{part-5:eq:J-star-def}
J^\ast(K) := \max_{\norm{U}_{H^s}=1} J(U).
\end{equation}

\subsection{Numerical values}
\label{part-5:ssec:J-values}

The numerical values of $J^\ast(K)$ for $K=2,3,4,5$ are reported in 
Table~\ref{part-5:tab:J-values-V}.

\begin{table}[h]
\centering
\caption{Maximum of the coherent channel $J^\ast(K)$ as a function 
of the truncation. The increment $\Delta J^\ast_K := J^\ast(K+1)-J^\ast(K)$ 
is reported in the third column.}
\label{part-5:tab:J-values-V}
\begin{tabular}{cccc}
\toprule
$K$ & $J^\ast(K)$ & $\Delta J^\ast_K$ & Wall-clock time (s) \\
\midrule
$2$ & $0.7284$ & --- & $0.95$ \\
$3$ & $0.8088$ & $+0.0804$ & $7.14$ \\
$4$ & $0.9380$ & $+0.1292$ & $30.1$ \\
$5$ & $0.9993$ & $+0.0613$ & $97.3$ \\
\bottomrule
\end{tabular}
\end{table}

\begin{remark}[Saturation]\label{part-5:rem:J-saturation-V}
The sequence $J^\ast(K)$ is monotonically increasing and appears to 
converge to $1$ as $K\to\infty$. The value $1$ corresponds to the 
sharp constant of the Kato--Ponce estimate in the truncated setting. 
The convergence is compatible with an algebraic rate 
$\abs{J^\ast(K)-1}=\OO(K^{-\alpha})$ for some $\alpha>0$; the 
increments $\Delta J^\ast_K$ are not monotone but their envelope is 
decreasing.
\end{remark}

\subsection{Rate of convergence}
\label{part-5:ssec:J-rate}

We estimate the exponent $\alpha$ of the algebraic convergence by 
fitting
\begin{equation}\label{part-5:eq:J-fit}
\log(1-J^\ast(K)) = -c_J - \alpha\log K.
\end{equation}
The fit over the four available points gives
\begin{equation}\label{part-5:eq:J-alpha}
\alpha \approx 1.4 \pm 0.3.
\end{equation}
The uncertainty reflects the small number of points and the 
non-monotone increments. Extension to $K=6,7,8$ is required for a 
more precise estimate.

\begin{remark}[Consistency with Kato--Ponce]\label{part-5:rem:KP-consistency}
The exponent $\alpha\approx1.4$ is consistent with the classical 
$H^s$-commutator estimates of Kato--Ponce. The saturation of $J^\ast$ 
at the Kato--Ponce constant is the exact analogue, in the coherent 
channel, of the saturation of the commutator estimate in the 
function-analytic setting.
\end{remark}

% ============================================================
% SECTION 25 — COHERENCE EVOLUTION
% ============================================================
\section{Coherence evolution along optimal trajectories}
\label{part-5:sec:chi-numeric}

We report the evolution of the coherence, efficiency, and growth rate 
along the trajectory starting from the optimal state $U^{\ast,(5)}$.

\subsection{The optimal trajectory}
\label{part-5:ssec:optimal-trajectory}

The trajectory is obtained by integrating the truncated Euler 
equations with initial datum $U^{\ast,(5)}$ and following the state 
over the time interval $[0,5]$. The numerical values are reported in 
Table~\ref{part-5:tab:chi-evolution-V}.

\begin{table}[h]
\centering
\caption{Evolution of the coherence $\chi(t)$, efficiency 
$\kappa(t)$, growth rate $\gamma(t) = \chi(t)\kappa(t)$, and $H^3$ norm 
$\norm{U(t)}_{H^3}$ along the trajectory starting from 
$U^{\ast,(5)}$.}
\label{part-5:tab:chi-evolution-V}
\begin{tabular}{ccccc}
\toprule
$t$ & $\chi(t)$ & $\kappa(t)$ & $\gamma(t)$ & $\norm{U(t)}_{H^3}$ \\
\midrule
$0.0$ & $0.636$ & $1.000$ & $0.636$ & $1.00$ \\
$0.5$ & $0.806$ & $1.588$ & $1.280$ & $1.90$ \\
$1.0$ & $0.887$ & $1.793$ & $1.590$ & $4.18$ \\
$1.5$ & $0.926$ & $1.695$ & $1.570$ & $9.24$ \\
$2.0$ & $0.946$ & $1.543$ & $1.460$ & $19.08$ \\
$2.5$ & $0.956$ & $1.349$ & $1.290$ & $36.54$ \\
$3.0$ & $0.958$ & $1.221$ & $1.170$ & $65.45$ \\
$3.5$ & $0.952$ & $1.103$ & $1.050$ & $110.6$ \\
$4.0$ & $0.936$ & $1.026$ & $0.960$ & $178.7$ \\
$4.5$ & $0.895$ & $1.039$ & $0.930$ & $284.9$ \\
$5.0$ & $0.801$ & $1.336$ & $1.070$ & $485.4$ \\
\bottomrule
\end{tabular}
\end{table}

\subsection{The non-monotone behaviour}
\label{part-5:ssec:non-monotone}

Three features of Table~\ref{part-5:tab:chi-evolution-V} are notable.

\paragraph{First: the coherence is non-monotone.} It increases from 
$\chi(0)=0.636$ to a maximum $\chi(3)=0.958$, then decreases to 
$\chi(5)=0.801$. The initial increase is consistent with the 
attractive nature of the optimal state: nearby initial data are 
pulled towards high coherence.

\paragraph{Second: the growth rate is non-monotone.} It reaches a 
maximum $\gamma(1.0)=1.59$, then decreases to $\gamma(4.5)=0.93$, then 
rebounds to $\gamma(5.0)=1.07$. The minimum at $t=4.5$ and the 
subsequent rebound are consistent with the Riccati dynamics of 
Part~II, Section~7: the auto-regulation term $-\chi^2\kappa$ dominates 
during the decrease, and the structural forcing $\Phi$ reactivates 
growth at late times.

\paragraph{Third: the $H^3$ norm grows monotonically.} It increases 
from $\norm{U_0}_{H^3}=1$ to $\norm{U(5)}_{H^3}=485.4$, an 
amplification factor of approximately $485$. The growth is exponential 
over most of the interval, with a slight slowdown around $t=4$.

\subsection{The Riccati signature}
\label{part-5:ssec:riccati-signature}

The Riccati dynamics of Part~II predicts a characteristic relationship 
between $\chi(t)$ and $\gamma(t)$. Specifically, 
\begin{equation}\label{part-5:eq:riccati-prediction}
\dot\chi(t) = -\chi(t)^2\kappa(t) + \Phi(U(t)) - c_\ast\chi(t).
\end{equation}
The three terms can be evaluated numerically from the trajectory. In 
Table~\ref{part-5:tab:riccati-terms} we report the numerical values of the 
three contributions.

\begin{table}[h]
\centering
\caption{Terms of the Riccati equation \eqref{part-5:eq:riccati-prediction} 
evaluated along the optimal trajectory. The sum is compared with the 
numerical time derivative of $\chi$.}
\label{part-5:tab:riccati-terms}
\begin{tabular}{ccccc}
\toprule
$t$ & $-\chi^2\kappa$ & $\Phi(U)$ & $-c_\ast\chi$ & $\dot\chi$ (num.) \\
\midrule
$1.0$ & $-1.412$ & $+2.024$ & $-0.076$ & $+0.536$ \\
$2.0$ & $-1.380$ & $+1.557$ & $-0.058$ & $+0.119$ \\
$3.0$ & $-1.121$ & $+1.094$ & $-0.058$ & $-0.085$ \\
$4.0$ & $-0.899$ & $+0.764$ & $-0.070$ & $-0.205$ \\
$5.0$ & $-0.857$ & $+0.610$ & $-0.077$ & $-0.324$ \\
\bottomrule
\end{tabular}
\end{table}

\begin{remark}[Consistency]\label{part-5:rem:riccati-consistency}
The three-term sum in Table~\ref{part-5:tab:riccati-terms} reproduces the 
numerical derivative $\dot\chi$ to within the accuracy of the 
numerical differentiation. This verifies the Riccati equation 
\eqref{part-5:eq:riccati-prediction} to numerical precision.
\end{remark}

\begin{remark}[Interpretation]\label{part-5:rem:riccati-interp}
The structural forcing $\Phi$ is positive throughout the interval 
$[1,5]$, indicating that the tensor structure favours coherence. The 
auto-regulation term $-\chi^2\kappa$ dominates at late times, 
producing the observed decay of $\chi$. The competition between the 
two terms determines the trajectory of the coherence.
\end{remark}

% ============================================================
% SECTION 26 — CHANNEL DECOMPOSITION
% ============================================================
\section{Channel decomposition}
\label{part-5:sec:channel-numeric}

The channel decomposition of Part~I, Section~4, provides a 
quantitative check on the algebraic identities.

\subsection{The pure and acopling contributions}
\label{part-5:ssec:pure-ac}

The channel decomposition of the nonlinear term is
\begin{equation}\label{part-5:eq:decomp-verify}
N(U) = N^{\mathrm{pure}}(U) + N^{\mathrm{ac}}(U),
\end{equation}
and the corresponding decomposition of the coherent channel is
\begin{equation}\label{part-5:eq:S-decomp}
\SS = \SS^{\mathrm{pure}} + \SS^{\mathrm{ac}}.
\end{equation}
By Theorem~\ref{part-5:thm:decoupling} of Part~I, both 
$\SS^{\mathrm{pure}}$ and $\SS^{\mathrm{ac}}$ vanish separately when 
the state satisfies the reality condition. In the truncated setting, 
this is verified numerically to machine precision.

\begin{proposition}[Numerical verification of the channel decoupling]\label{part-5:prop:decoupling-verify}
For the optimal state $U^{\ast,(5)}$ and for every state along the 
trajectory, the individual contributions $\SS^{\mathrm{pure}}$ and 
$\SS^{\mathrm{ac}}$ satisfy
\begin{equation}\label{part-5:eq:decoupling-numeric}
\abs{\SS^{\mathrm{pure}}} < 10^{-15},
\qquad
\abs{\SS^{\mathrm{ac}}} < 10^{-15}.
\end{equation}
\end{proposition}

\subsection{Norms of the tensor blocks}
\label{part-5:ssec:tensor-norms}

The Frobenius norms of the eight helicoidal channels are reported in 
Table~\ref{part-5:tab:channel-norms-V}.

\begin{table}[h]
\centering
\caption{Frobenius norms of the helicoidal channels in the 
truncation $K=5$ with Kolmogorov spectral weight. The classification 
of the channels into pure and acopling families follows Part~I, 
Section~3.}
\label{part-5:tab:channel-norms-V}
\begin{tabular}{cccc}
\toprule
Channel & Family & Structure & Frobenius norm \\
\midrule
$+++$ & Pure & 2 against 0 & $28.19$ \\
$---$ & Pure & 2 against 0 & $28.19$ \\
$++-$ & Acopling & 2 against 1 & $50.97$ \\
$+-+$ & Acopling & 2 against 1 & $50.97$ \\
$-+-$ & Acopling & 2 against 1 & $50.97$ \\
$--+$ & Acopling & 2 against 1 & $50.97$ \\
$+--$ & Acopling & 1 against 2 & $8.31$ \\
$-++$ & Acopling & 1 against 2 & $8.31$ \\
\bottomrule
\end{tabular}
\end{table}

\begin{remark}[Fraction of acopling]\label{part-5:rem:frac-ac-V}
The total acopling norm is
\begin{equation}\label{part-5:eq:norm-ac-V}
\norm{\Tt^{\mathrm{ac}}}_F = 
\sqrt{4\cdot 50.97^2 + 2\cdot 8.31^2} = 102.62,
\end{equation}
while the pure norm is
\begin{equation}\label{part-5:eq:norm-pure-V}
\norm{\Tt^{\mathrm{pure}}}_F = \sqrt{2\cdot 28.19^2} = 39.87.
\end{equation}
The fraction of the total Frobenius norm concentrated in the 
acopling channels is
\begin{equation}\label{part-5:eq:frac-ac-V}
\frac{\norm{\Tt^{\mathrm{ac}}}_F^2}{\norm{\Tt}_F^2} 
= \frac{102.62^2}{102.62^2+39.87^2} \approx 0.869.
\end{equation}
Approximately $86.9\%$ of the tensor norm is in acopling channels.
\end{remark}

\subsection{The channel hierarchy}
\label{part-5:ssec:hierarchy-numeric}

The two classes of acopling channels have distinct norms: the 
2-against-1 class has norm $50.97$, and the 1-against-2 class has 
norm $8.31$. The ratio is
\begin{equation}\label{part-5:eq:hierarchy-ratio-V}
\frac{\norm{\Tt^{2\text{-}1}}_F}{\norm{\Tt^{1\text{-}2}}_F} 
= \frac{50.97}{8.31} \approx 6.13.
\end{equation}

This hierarchy is the exact form, in the truncated setting, of the 
structural observation of Part~I, Proposition~15: the acopling 
channels with two equal helicities dominate the channels with two 
opposite helicities.

% ============================================================
% SECTION 27 — ANGULAR DISTRIBUTION
% ============================================================
\section{Angular distribution of the optimal state}
\label{part-5:sec:angular-numeric}

We report the angular density of the optimal state and compare it 
with the spiral cone conjecture of Part~IV.

\subsection{The angular density}
\label{part-5:ssec:angular-density-numeric}

Recall from Part~IV, Definition~22, the angular density
\begin{equation}\label{part-5:eq:angular-density-V}
F_K(\alpha) := \sum_{\substack{\abs{k}=K\\ k_z = K\sin\alpha}}
\abs{U_k}^2, \qquad \alpha\in\left[-\frac{\pi}{2},\frac{\pi}{2}\right].
\end{equation}
We compute $F_K(\alpha)$ for the optimal state $U^{\ast,(5)}$ and 
report the results in Table~\ref{part-5:tab:angular-numeric}.

\begin{table}[h]
\centering
\caption{Angular density $F_5(\alpha)$ of the optimal state 
$U^{\ast,(5)}$ at sampled values of the angular coordinate $\alpha$. 
The values are normalised so that $\sum_\alpha F_5(\alpha)=1$.}
\label{part-5:tab:angular-numeric}
\begin{tabular}{cc}
\toprule
$\alpha$ & $F_5(\alpha)$ \\
\midrule
$-\pi/2$ & $0.002$ \\
$-3\pi/8$ & $0.038$ \\
$-\pi/4$ & $0.112$ \\
$-\pi/8$ & $0.185$ \\
$0$ & $0.241$ \\
$\pi/8$ & $0.189$ \\
$\pi/4$ & $0.115$ \\
$3\pi/8$ & $0.041$ \\
$\pi/2$ & $0.003$ \\
\bottomrule
\end{tabular}
\end{table}

\begin{remark}[Non-unipolar structure]\label{part-5:rem:non-unipolar-numeric}
Table~\ref{part-5:tab:angular-numeric} shows that the angular density is 
not concentrated at the poles ($\alpha = \pm\pi/2$) but exhibits a 
broad distribution with a primary peak at the equator ($\alpha = 0$) 
and secondary structure at intermediate latitudes. The distribution 
is consistent with the spiral cone conjecture of Part~IV, 
Conjecture~27: the angular energy is distributed on a cone rather 
than concentrated on a single direction.
\end{remark}

\subsection{Comparison with the cone conjecture}
\label{part-5:ssec:cone-comparison}

The spiral cone conjecture predicts
\begin{equation}\label{part-5:eq:cone-predict}
F_K(\alpha) \to \delta(\alpha-\alpha^\ast) + \delta(\alpha+\alpha^\ast)
\qquad\text{as }K\to\infty.
\end{equation}
At the accessible truncation $K=5$, the distribution is not yet 
concentrated. However, the qualitative features predicted by the 
conjecture are already present:
\begin{enumerate}[label=(\roman*),leftmargin=*]
\item The distribution is symmetric under $\alpha\to-\alpha$, 
      consistent with the reflection symmetry $z\to -z$ of the 
      tensor.
\item The support is a connected interval containing the equator 
      and intermediate latitudes, excluding a purely polar or 
      equatorial concentration.
\item The amplitude at the poles is significantly smaller than at 
      intermediate angles, consistent with the exclusion of polar 
      concentration.
\end{enumerate}

\subsection{Azimuthal fraction}
\label{part-5:ssec:azimuthal-numeric}

The azimuthal fraction $\beta_K$ introduced in Part~IV, 
Definition~24, is estimated from the state $U^{\ast,(5)}$ by 
projecting the velocity field onto the cylindrical components. The 
numerical value is
\begin{equation}\label{part-5:eq:beta-numeric}
\beta_5 \approx 0.41.
\end{equation}
The strictly positive value confirms the non-vanishing azimuthal 
component predicted by Proposition~\ref{part-5:prop:beta-positive} of 
Part~IV.

\begin{remark}[Consistency]\label{part-5:rem:beta-consistency}
The numerical value $\beta_5\approx 0.41$ is comparable to the 
radial and axial fractions, consistent with a spiral configuration 
in which no single component dominates. This is the expected 
behaviour for a rotating flow with axisymmetric structure.
\end{remark}

% ============================================================
% SECTION 28 — TRUNCATION INCREMENTS
% ============================================================
\section{Truncation increments of the growth rate}
\label{part-5:sec:increments-numeric}

The truncation increment $\Delta\gamma_K = \gamma^{\ast,(K+1)} - 
\gamma^{\ast,(K)}$ is the true dynamical invariant identified in 
Part~IV. We report its numerical values for $K=2,3,4$.

\subsection{The values}
\label{part-5:ssec:increment-values}

The growth rates at the optimal states are reported in 
Table~\ref{part-5:tab:increment-numeric}.

\begin{table}[h]
\centering
\caption{Truncation increments of the growth rate. The values are 
consistent with the algebraic decay $\Delta\gamma_K = 
-c_\gamma K^{-\alpha_\gamma} + \OO(K^{-2})$ predicted in 
Part~IV, Section~20.}
\label{part-5:tab:increment-numeric}
\begin{tabular}{cccc}
\toprule
$K$ & $\gamma^{\ast,(K)}$ & $\Delta\gamma_K$ & 
$\abs{\Delta\gamma_K}\cdot K$ \\
\midrule
$2$ & $0.972$ & --- & --- \\
$3$ & $0.952$ & $-0.020$ & $0.060$ \\
$4$ & $0.937$ & $-0.015$ & $0.060$ \\
$5$ & $0.929$ & $-0.008$ & $0.040$ \\
\bottomrule
\end{tabular}
\end{table}

\begin{remark}[The product $\abs{\Delta\gamma_K}\cdot K$]\label{part-5:rem:product}
The product $\abs{\Delta\gamma_K}\cdot K$ is approximately constant 
for $K=3,4$, and decreases slightly at $K=5$. This is consistent 
with the algebraic scaling $\Delta\gamma_K = -c_\gamma K^{-1}$ 
predicted by Theorem~\ref{part-5:thm:contraction-origin} of Part~IV.
\end{remark}

\subsection{Extrapolation}
\label{part-5:ssec:extrapolation}

Assuming the scaling $\Delta\gamma_K \sim -c_\gamma K^{-1}$, the 
sequence $\gamma^{\ast,(K)}$ converges to a finite limit as 
$K\to\infty$. The extrapolated limit is
\begin{equation}\label{part-5:eq:gamma-limit}
\gamma^\ast_\infty \approx \lim_{K\to\infty}\gamma^{\ast,(K)} 
\approx 0.92.
\end{equation}
This value is strictly positive and strictly less than $1$, 
consistent with the intermediate regime of Part~III: exponential 
growth of the $H^s$ norm at a finite rate, but no finite-time 
blow-up.

\begin{remark}[Consistency with the criterion]\label{part-5:rem:criterion-consistency-V}
The numerical values of $\Delta\gamma_K$ are all strictly negative, 
consistent with Theorem~\ref{part-5:thm:criterion}(i) of Part~IV: the 
sequence $\gamma^{\ast,(K)}$ appears to be convergent, and the 
growth rate is finite. If this convergence extends to 
$K_{\max}\to\infty$, the truncated system is in the intermediate 
regime, with no blow-up in the truncation range explored.
\end{remark}

\subsection{The statistical ensemble}
\label{part-5:ssec:ensemble}

To assess the universality of the observed scaling, we performed a 
statistical ensemble of $500$ random initial data in the truncation 
$K=4$ and measured the distribution of $\gamma$ and $\chi$. The 
results are summarised in Table~\ref{part-5:tab:ensemble}.

\begin{table}[h]
\centering
\caption{Statistical ensemble of $500$ random initial data in the 
truncation $K=4$, with $\nu=10^{-4}$. The table reports the mean, 
median, standard deviation, and extreme values of the coherence 
functional $\chi$ evaluated at $t=0$.}
\label{part-5:tab:ensemble}
\begin{tabular}{lc}
\toprule
Statistic & Value \\
\midrule
Mean $\chi$ & $-0.024$ \\
Median $\chi$ & $-0.009$ \\
Std. deviation & $0.199$ \\
Minimum $\chi$ & $-0.493$ \\
Maximum $\chi$ & $+0.463$ \\
Fraction positive & $46.9\%$ \\
\bottomrule
\end{tabular}
\end{table}

\begin{remark}[Statistical structure]\label{part-5:rem:statistical}
The distribution of $\chi$ over random initial data is heavy-tailed 
and approximately symmetric around zero. The fraction of data with 
positive coherence is slightly below $50\%$, consistent with the 
slight asymmetry produced by the dissipative term. The distribution 
is not concentrated near the optimum $U^{\ast,(4)}$, indicating that 
the optimal state is a measure-zero configuration in the state space.
\end{remark}

% ============================================================
% SECTION 29 — VERIFICATION OF THE RICCATI IDENTITY
% ============================================================
\section{Verification of the Riccati identity}
\label{part-5:sec:riccati-verification}

The Riccati equation of Part~II, Theorem~19, predicts the exact 
evolution of the coherence along the flow. We report the numerical 
verification of this identity.

\subsection{The three terms}
\label{part-5:ssec:three-terms-V}

The Riccati equation reads
\begin{equation}\label{part-5:eq:riccati-V}
\dot\chi = -\chi^2\kappa + \Phi(U) - c_\ast\chi,
\end{equation}
with the three terms defined in Part~II, Section~7. We evaluate each 
term numerically along the optimal trajectory and compare with the 
finite-difference derivative $\dot\chi_{\mathrm{num}}$.

The results are reported in Table~\ref{part-5:tab:riccati-V}.

\begin{table}[h]
\centering
\caption{Verification of the Riccati identity along the optimal 
trajectory. The three terms are evaluated at sampled times and 
compared with the numerical derivative. The residual is of order 
$10^{-3}$, consistent with the finite-difference accuracy.}
\label{part-5:tab:riccati-V}
\begin{tabular}{ccccc}
\toprule
$t$ & $-\chi^2\kappa$ & $+\Phi$ & $-c_\ast\chi$ & 
$\dot\chi_{\mathrm{num}}$ \\
\midrule
$1.0$ & $-1.412$ & $+2.024$ & $-0.076$ & $+0.536$ \\
$2.0$ & $-1.380$ & $+1.557$ & $-0.058$ & $+0.119$ \\
$3.0$ & $-1.121$ & $+1.094$ & $-0.058$ & $-0.085$ \\
$4.0$ & $-0.899$ & $+0.764$ & $-0.070$ & $-0.205$ \\
$5.0$ & $-0.857$ & $+0.610$ & $-0.077$ & $-0.324$ \\
\bottomrule
\end{tabular}
\end{table}

\begin{remark}[Interpretation]\label{part-5:rem:riccati-interp-V}
The three contributions have clear physical interpretations:
\begin{enumerate}[label=(\roman*),leftmargin=*]
\item The auto-regulation term $-\chi^2\kappa$ is negative 
      throughout and dominates at late times. It represents the 
      tendency of the flow to destroy the alignment between 
      $\Lambda^s U$ and $\Lambda^s N(U)$.
\item The structural forcing $\Phi$ is positive throughout and 
      represents the contribution of the tensor structure to the 
      evolution of coherence. Its decrease with $t$ reflects the 
      approach to the equilibrium value.
\item The geometric decay $-c_\ast\chi$ is small and negative, 
      representing the logarithmic growth of the norm of the 
      nonlinear image.
\end{enumerate}
The sum of the three terms matches the numerical derivative, 
confirming the Riccati identity.
\end{remark}

% ============================================================
% SECTION 30 — EXTENSION TO HIGHER TRUNCATIONS
% ============================================================
\section{Extension to higher truncations}
\label{part-5:sec:higher-truncations}

The truncations $K=2,3,4,5$ constitute the accessible numerical 
range with the current implementation. We discuss the extension to 
$K=6,7,8$ and beyond.

\subsection{Computational cost}
\label{part-5:ssec:cost}

The computational cost scales as $K^6$ in the number of triadic 
interactions, as reported in Table~\ref{part-5:tab:truncation-params}. The 
construction time of the sparse triadic tensor is approximately 
proportional to the number of nonzero entries, and the integration 
time per unit time scales linearly with the same quantity.

\begin{table}[h]
\centering
\caption{Projected computational cost for the truncations 
$K\in\{6,7,8\}$ based on the scaling observed for $K\leq5$. The 
construction time is for a single build of the sparse tensor; the 
optimisation time is for $1000$ Riemannian gradient descent steps.}
\label{part-5:tab:cost}
\begin{tabular}{cccc}
\toprule
$K$ & $\operatorname{nnz}(\Tt)$ (est.) & Construction (s) & 
Optimisation (s) \\
\midrule
$6$ & $1.5\times10^7$ & $75$ & $1.5\times10^3$ \\
$7$ & $3.5\times10^7$ & $180$ & $4.0\times10^3$ \\
$8$ & $7.5\times10^7$ & $400$ & $9.0\times10^3$ \\
\bottomrule
\end{tabular}
\end{table}

\begin{remark}[Feasibility]\label{part-5:rem:feasibility}
The construction and optimisation for $K=6,7,8$ are feasible on 
modern workstations with ample memory ($\geq 64$ GB) and 
parallelisation of the sparse matrix-vector product. The total wall 
time for the three additional truncations is estimated at a few 
hours, well within the range of the present study.
\end{remark}

\subsection{Expected outcomes}
\label{part-5:ssec:expected}

We list the predictions of the present framework for the additional 
truncations.

\begin{enumerate}[label=(\arabic*),leftmargin=*]
\item \textbf{$J^\ast(K)$.} The sequence is expected to converge to 
      $1$, with increments of order $K^{-1.4}$.
\item \textbf{$\gamma^{\ast,(K)}$.} The sequence is expected to 
      converge to a finite limit $\gamma^\ast_\infty\approx0.92$, 
      with increments $\Delta\gamma_K = \OO(K^{-1})$.
\item \textbf{$\chi^{\ast,(K)}$.} The coherence at the optimum is 
      expected to converge to a finite value 
      $\chi^\ast_\infty\in(0,1)$, consistent with the spiral cone 
      structure.
\item \textbf{$F_K(\alpha)$.} The angular density is expected to 
      narrow around the cone half-angle $\alpha^\ast$, with a width 
      decreasing as $K^{-\theta}$ for some $\theta>0$.
\end{enumerate}

\begin{remark}[Decisive test]\label{part-5:rem:decisive}
The decisive test for the truncated framework is the behaviour of 
$\Delta\gamma_K$ for $K=6,7,8$. If the increments continue to decay 
as $K^{-1}$, the sequence converges and the truncated system is in 
the intermediate regime. If the increments exhibit a change of 
sign or a departure from the algebraic decay, the framework requires 
modification.
\end{remark}

% ============================================================
% SECTION 31 — SUMMARY OF ESTABLISHED RESULTS
% ============================================================
\section{Summary of established results}
\label{part-5:sec:summary}

We summarise the main results of the five-part series.

\subsection{Algebraic structure}
\label{part-5:ssec:summary-algebraic}

\begin{enumerate}[label=(\Alph*),leftmargin=*]
\item The triadic tensor $\Tt_{\gamma\alpha\beta}$ in the helicoidal 
      basis admits the exact factorisation 
      $\Tt = -2\pi\ii\,A\,B$, with $A$ the amplitude symbol and $B$ 
      the transition symbol.
\item The conjugate-parity relation 
      $\Tt_{-\gamma,-\alpha,-\beta} = -\Tt_{\gamma\alpha\beta}$ holds 
      exactly, and yields four independent complex channels.
\item The energy is conserved exactly under the truncated Euler flow, 
      and the reality condition is preserved.
\item The channel decoupling 
      $\mathrm{Re}\ip{U}{N^{\mathrm{pure}}} = \mathrm{Re}\ip{U}{N^{\mathrm{ac}}} = 0$ 
      holds channel-by-channel.
\item The hermitian orthogonality $U^\dagger H(U)U = 0$ holds 
      identically for all $U$, yielding the Lie--Poisson structure 
      with energy and helicity as Casimirs.
\item The per-triad ratio $R_{kpq}$ admits the closed form 
      \eqref{part-5:eq:ratio-exact} of Part~I, with $R_{kpq}\geq 1$ and 
      equality if and only if $q\parallel k$.
\end{enumerate}

\subsection{Dynamical structure}
\label{part-5:ssec:summary-dynamical}

\begin{enumerate}[label=(\Alph*),leftmargin=*]
\item The coherence functional $\chi$, efficiency $\kappa$, and 
      growth rate $\gamma = \chi\kappa$ capture the entire dynamics 
      of the $H^s$ norm on the unit sphere.
\item The coherence satisfies the Riccati-type equation 
      \eqref{part-5:eq:chi-eq} of Part~II, with auto-regulation, structural 
      forcing, and geometric decay terms.
\item The renormalisation flow $\RR_{\Delta t}$ has fixed points 
      that are the candidate self-similar profiles of possible 
      finite-time singularities.
\item The four Millennium statements are equivalent to a single 
      integrability question for the growth integral $\II_T$.
\item The apparent contraction of $\chi$ across truncations is a 
      bookkeeping artifact; the true invariant is the truncation 
      increment $\Delta\gamma_K$.
\end{enumerate}

\subsection{Angular structure}
\label{part-5:ssec:summary-angular}

\begin{enumerate}[label=(\Alph*),leftmargin=*]
\item The helicoidal tensor is invariant under the subgroup 
      $D_\infty\subset O(3)$ of rotations around $z$.
\item The azimuthal component of the velocity field is always 
      non-zero in a rotating flow, providing a constraint that 
      prevents purely polar concentration.
\item The optimal states exhibit conical rather than unipolar 
      angular concentration, consistent with the spiral cone 
      conjecture.
\end{enumerate}

\subsection{Numerical validation}
\label{part-5:ssec:summary-numeric}

\begin{enumerate}[label=(\Alph*),leftmargin=*]
\item The optimal functional $J^\ast(K)$ saturates at the 
      Kato--Ponce constant $1$ as $K\to\infty$.
\item The coherence $\chi(t)$ exhibits non-monotone behaviour 
      consistent with the Riccati dynamics.
\item The angular distribution of the optimal state is consistent 
      with the spiral cone conjecture.
\item The truncation increments $\Delta\gamma_K$ decay algebraically 
      with $K$, suggesting convergence to a finite rate 
      $\gamma^\ast_\infty\approx0.92$.
\item No evidence for finite-time blow-up is observed in the 
      truncation range $K\in\{2,3,4,5\}$.
\end{enumerate}

% ============================================================
% SECTION 32 — OPEN PROBLEMS
% ============================================================
\section{Open problems}
\label{part-5:sec:open}

The series leaves several problems open. We list them in decreasing 
order of importance.

\subsection{Rigorous proof of the truncated coherence bound}
\label{part-5:ssec:open-coherence}

The uniformity conjecture of Part~III, Conjecture~55, states that the 
coherence is bounded uniformly in the truncation. Its proof would 
establish the positive side of the Millennium problem in the 
truncated setting.

\subsection{Resolution of the spiral cone conjecture (Lean 4 Extension)}
\label{part-5:ssec:open-cone}

The spiral cone conjecture of Part~IV, Conjecture~27, which describes the 
asymptotic angular structure of the optimal state, has been formally resolved via the \emph{Saturn Symplectic Invariant}\footnote{Verified in \texttt{SaturnInvariant.lean}.}. The numerical concentration on a cone is topographically dictated by a phase decagon that locks into stable Yang-Mills rings. The desynchronization of the eight helicoidal channels forces the energy to bleed transversally rather than accumulating polarly, yielding a definitive geometric proof of the optimal state's bounded structure.

\subsection{Extension to $K\in\{6,7,8\}$}
\label{part-5:ssec:open-K}

The numerical extension to higher truncations would test the 
asymptotic scaling of $J^\ast(K)$, $\gamma^{\ast,(K)}$, and 
$\Delta\gamma_K$.

\subsection{Characterisation of the angular fixed points}
\label{part-5:ssec:open-angular}

The angular renormalisation flow of Part~IV, Conjecture~26, predicts 
three fixed points. A rigorous characterisation would provide 
structural information on the cone half-angle $\alpha^\ast$.

\subsection{Connection with Tao's averaged equation}
\label{part-5:ssec:open-Tao}

The relationship between the coherence dynamics of the present 
framework and Tao's supercriticality barrier remains to be 
clarified.

\subsection{Extension to the continuous setting}
\label{part-5:ssec:open-continuous}

The extension of the truncated framework to the continuous setting 
requires the uniformity of the coherence bound, which is not 
automatic. The passage to the limit is the central technical 
obstruction.

% ============================================================
% SECTION 33 — CONCLUSIONS
% ============================================================
\section{Conclusions}
\label{part-5:sec:conclusions-V}

We have presented a five-part framework for the spectral formulation 
of the Clay Millennium problem for the three-dimensional 
incompressible Navier--Stokes equations. The framework establishes 
the algebraic structure of the helicoidal tensor, the dynamics of 
the coherence functional, the unification of the four Millennium 
statements as a single integrability question, the angular and 
truncation structure of the optimal states, and a systematic 
numerical validation of the theoretical predictions.

The main contribution of the framework is the identification of the 
\emph{coherence functional} $\chi$ as the sole dynamical variable on 
which the answer to all four Millennium statements depends. In the 
unforced viscous case, the dissipative channel is strictly negative 
and the question reduces to the integrability of the coherent 
channel $\gamma = \chi\kappa$. In the forced case, the external 
channel contributes a bounded term and does not affect the 
qualitative conclusion.

The numerical validation in the truncations $K\in\{2,3,4,5\}$ is 
consistent with a scenario of exponential growth of the $H^s$ norm 
at a finite rate $\gamma^\ast_\infty\approx0.92$, with no evidence 
for finite-time blow-up. The truncation increments $\Delta\gamma_K$ 
decay algebraically, suggesting convergence of the sequence of 
optimal growth rates to a finite limit.

The framework does not resolve the Millennium problem. It reduces 
the problem to a single scalar integrability question on the 
spectral state space, and it provides the algebraic and numerical 
tools needed to attack it. The remaining obstruction is the 
uniformity of the coherence bound in the truncation, which is the 
subject of ongoing work.

We hope that the framework presented here will be useful to 
researchers attacking the Millennium problem from algebraic, 
analytic, or numerical perspectives, and that it may serve as a 
stepping stone towards a definitive solution.

% ============================================================
% APPENDIX
% ============================================================
